\documentclass[11pt]{article}
\usepackage[margin=1in]{geometry}
\usepackage{amsmath,amssymb,amsthm,mathtools,booktabs,longtable,array,microtype}
\usepackage[T1]{fontenc}
\usepackage{lmodern}
\usepackage{xcolor,listings,enumitem,fancyhdr,xurl,needspace}
\usepackage[hidelinks]{hyperref}
% Generated by code/build_document_data.py; do not edit by hand.
\newcommand{\Nstates}{222}
\newcommand{\LastState}{221}
\newcommand{\Saving}{16}
\newcommand{\Nrows}{444}
\newcommand{\Nsites}{3891}
\newcommand{\Ndag}{204}
\newcommand{\Ndepth}{45}
\newcommand{\Nzeros}{6}
\newcommand{\Ncases}{5060}
\newcommand{\BootSteps}{95714}
\newcommand{\ClauseOrder}{lower-first}
\newcommand{\Hbody}{I(L_n,I(U_n,F))}
\newcommand{\DenseHash}{f3c8492394b2589c47cfe0bb38a3dc6e51e11bfdead7cdec6b07913a4106cb3e}
\newcommand{\NativeHash}{88a9c5ea637560bf58dba612d598663213884d7bb7981230164bc058babaa427}
\newcommand{\SourceHash}{1d258cfe56f8b845826a63af57de8134fbc4886b47f95312cf6d49424f3f3072}

\hypersetup{pdftitle={A 222-state binary Turing machine for unary power covering},
 pdfauthor={Research prepared for Vladimir Reshetnikov},
 pdfsubject={A PA-provable consistency equivalence and exact machine verification},
 pdfkeywords={Turing machine, PA, consistency, unary power covering, U222}}
\newcommand{\PA}{\mathsf{PA}}
\newcommand{\ThU}{\mathsf U}
\newcommand{\ThW}{\mathsf W}
\newcommand{\IR}{\mathsf{IR}}
\newcommand{\Con}{\operatorname{Con}}
\newcommand{\Halt}{\operatorname{Halt}}
\newcommand{\Sub}{\operatorname{Sub}}
\newcommand{\Emp}{\operatorname{Emp}}
\newcommand{\Num}{\operatorname{Num}}
\newcommand{\Succ}{\operatorname{Succ}}
\newcommand{\Eq}{\mathrel{\simeq}}
\newcommand{\CU}{\mathcal C_U}
\newcommand{\Nn}{\mathbb N}
\newtheorem{theorem}{Theorem}[section]
\newtheorem{lemma}[theorem]{Lemma}
\newtheorem{proposition}[theorem]{Proposition}
\newtheorem{corollary}[theorem]{Corollary}
\theoremstyle{definition}\newtheorem{definition}[theorem]{Definition}
\pagestyle{fancy}\fancyhf{}\fancyhead[L]{\small Unary power covering and a small consistency machine}
\fancyhead[R]{\small 28 October 2026}\fancyfoot[C]{\thepage}
\setlength{\headheight}{14pt}
\setlength{\emergencystretch}{3em}
\lstset{basicstyle=\ttfamily\small,breaklines=true,columns=fullflexible,
  frame=single,rulecolor=\color{black!20},xleftmargin=0.5em,xrightmargin=0.5em,
  aboveskip=1em,belowskip=1em,showstringspaces=false}
\title{A \Nstates-state binary Turing machine\texorpdfstring{\\}{ }for the consistency of unary power covering}
\author{Research prepared for Vladimir Reshetnikov}
\date{28 October 2026}
\begin{document}
\maketitle
\begin{abstract}
We construct an explicit deterministic, single-tape, two-symbol Turing machine
with \Nstates\ working states, including initialization, whose blank-tape
non-halting is provably equivalent in Peano arithmetic to the syntactic
consistency of the supplied unary power-covering theory~$\ThU$.
This improves the requested 238-state benchmark by \Saving\ states.
The construction replaces the Theory~K axiom emitter in the accepted K215
proof-search architecture by a fixed-variable power-covering emitter. The
matrix code itself supplies a fresh output-variable label; both exceptional
small codes are proved harmless. The retained logical
calculus and the new schema are connected by uniform finite proof
transformations. The literal table is checked against every one of its
\Ncases\ primitive simulation cases using arbitrary natural register values,
an arbitrary gap, and an arbitrary nonempty finite spare list. Initialization is
independently replayed from a genuinely blank tape.

The result is an ordinary mathematical proof, with a PA formalizability
argument and executable symbolic verification. This release does not contain
a newly compiled Lean, Rocq, or MM0 certificate for its final theorem. No
absolute state minimality or unconditional consistency assertion is made.
\end{abstract}

\section{Result, conventions, and scope}

\begin{theorem}[Main result]\label{thm:main}
Let $M_U$ be the exact transition table in \texttt{machine/u.tm}.
It has \Nstates\ working states and satisfies
\[
 \PA\vdash \neg\Halt(M_U,\mathrm{blank})\ \longleftrightarrow\ \Con(\ThU).
\]
Using the uniform interpretation theorem in the supplied $\ThU$ report,
\[
 \PA\vdash \neg\Halt(M_U,\mathrm{blank})\ \longleftrightarrow\ \Con(\PA).
\]
\end{theorem}

Here $\Con(T)$ is the assertion that no natural number codes a finite
classical first-order $T$-refutation. The assertion on the left says that no
natural number codes a finite halting run of the displayed machine.
The proof concerns these arithmetized assertions, not just their truth in
the standard natural numbers. In particular, all simulation and proof
translations below are justified by induction available in PA.

The tape is indexed by $\mathbb Z$, starts entirely with symbol~0, and has
one head initially at position~0. Every step writes one symbol from
$\{0,1\}$, moves exactly one cell left or right, and changes state. The
dense presentation starts in state~0, has working labels
$0,\ldots,\LastState$, and has a distinct halt label~\Nstates\ with no
outgoing transitions. Every working state has both binary transition rows.
The count is therefore \Nstates, not \Nstates\ plus a hidden initialization
program or an externally prepared input tape. The native presentation uses
halt label~254 and dispatcher~84; the included bijection translates every
transition and the start state exactly.

\begin{center}
\begin{tabular}{lr}
\toprule
Quantity & Exact value\\\midrule
Working states & \Nstates\\
Transition rows & \Nrows\\
Source primitive occurrences & \Nsites\\
Distinct source subtrees, including primitive leaves & \Ndag\\
Maximum source depth / consecutive zeros & \Ndepth\ / \Nzeros\\
Zero/positive primitive cases checked & \Ncases\\
Literal blank initialization steps & \BootSteps\\
\bottomrule
\end{tabular}
\end{center}

\paragraph{What is and is not inherited.}
The K215 project supplies an exceptionally compact register interpreter,
proof-search source, and finite logical proof templates
\cite{proveit,k215pa}. Its ZFC endpoint is a theorem about a different
transition table and a different nonlogical theory. We do not apply that
endpoint to the modified table. We establish the new emitter contract,
the new theory calibration, the exact modified source contract, and the
complete literal-machine simulation separately.

The 238-state PA reference is the construction described in the upstream
\texttt{PA/m238} development \cite{m238}. Our comparison uses the same
two-symbol blank-tape convention. We call \Nstates\ the best verified
construction in this release; we do not prove that every smaller machine
fails the requested property, or assert a globally exhaustive priority search.

\section{The theory and the useful machine-facing extension}

The supplied theory $\ThU$ uses classical equality-free first-order logic
over a nonempty domain and one binary relation $E$, written as membership.
Set
\[
 \Sub(x,a):=\forall t\,(E(t,x)\to E(t,a)).
\]
For every formula $\varphi(x)$ with no other free variable, its sole
nonlogical schema is the universal closure of
\begin{equation}\label{eq:U}
 \forall a\,\exists b\left[
  \forall x\,(E(x,b)\to\varphi(x))\ \land\
  \forall x\,(\Sub(x,a)\to(\varphi(x)\to E(x,b)))
 \right].
\end{equation}
Arbitrarily many bound variables are allowed. The source variable $a$
cannot occur freely in a legal matrix of $\ThU$. The two clauses describe
a weak cover; they do not require every member of $b$ to be a subset of $a$.

Let $\ThW$ allow arbitrary additional matrix parameters, including the source,
in the same schema. Write
\[
 \IR:=\forall x\,\neg E(x,x).
\]
The machine can emit the unrestricted schema: recognizing the unary
free-variable restriction would add implementation work without changing
the consistency strength needed here.

\begin{proposition}[Uniform consistency reductions]\label{prop:theories}
There are PA-verifiable transformations of finite refutations yielding
\[
 \PA\vdash\Con(\ThU)\leftrightarrow\Con(\PA)
                 \leftrightarrow\Con(\ThW+\IR).
\]
The upper interpretation remains valid when ordinary equality and matrices
containing equality are permitted in the machine-facing presentation.
\end{proposition}
\begin{proof}
The first equivalence is the interpretation theorem proved in the supplied
report \cite{unary}, included in full under \texttt{references/theory\_U}.
Its relevant uniformity is summarized below. Since $\ThW$ includes every
$\ThU$ instance, a $\ThU$-refutation is immediately a
$\ThW+\IR$-refutation.

For the upper interpretation, interpret $E(x,a)$ as ``bit $x$ of the
binary expansion of $a$ is~1,'' and equality as arithmetic equality. PA
proves $E(x,a)\to x<a$ and hence irreflexivity. It also proves
$\Sub(x,a)\to x\le a$. For each fixed translated matrix
$\theta(x,a,\bar p)$, full PA induction constructs a number $b$ whose
1-bits are precisely the integers $x\le a$ satisfying
$\Sub(x,a)\land\theta(x,a,\bar p)$. Thus $b$ satisfies both weak-cover
clauses. This bounded finite selection uses the induction instance for the
particular formula $\theta$; it does not assume a uniform arithmetic truth
predicate. Arithmetic coding of exponentiation and finite bit sequences is
available in PA by its usual primitive-recursive definitional extensions.

Replacing each axiom in a $\ThW+\IR$ refutation by this finite arithmetic
derivation produces a PA-refutation. The $\PA$-to-$\ThU$ interpretation
then produces a $\ThU$-refutation. The translations are recursive on
formulas and finite proof lists, and PA verifies their correctness by
induction. The same upper argument handles equality-containing matrices.
\end{proof}

\subsection{The lower interpretation used from the supplied report}

For clarity about this dependency, the lower interpretation is much stronger
than an observed resemblance to finite sets. It uses coextensionality
$u\Eq v$ defined by
$\forall z(E(z,u)\leftrightarrow E(z,v))$, and numerical representatives
that are saturated, transitive, and linearly ordered by membership modulo
$\Eq$. Its endpoint condition requires each nonempty intersection of an
existing object's member profile with a numerical initial segment to have
least and greatest member classes. No external finiteness or well-foundedness
is assumed.

The dependency order is essential. Empty objects and local predecessor
facts are derived first. A unary filter of failures gives induction on
numerical representatives before total successor existence is known.
Parameters in an induction formula are internalized by a universally
quantified implication whose antecedent states that the corresponding base
and step premises hold. Induction then yields comparison and least failed
numerical tuples. Such least tuples permit bounded numerical comprehension
using genuinely unary matrices. For example, the matrix says that $z$
belongs to the definable profile associated with some least failing tuple;
the tuple is bound inside the matrix rather than used as an illicit free
parameter.

This comprehension theorem yields saturated successors, powers of numerical
objects, and saturated Kuratowski pair codes. Graphs with numerical
parameters are formed inside a suitable numerical power bound.
Deterministic histories are extended by a canonical formula that existentially
quantifies over prior histories; an arbitrary old graph is not inserted as
a free filter parameter. Successor iteration gives addition, and iteration
of addition gives multiplication. These operations are total and functional
modulo $\Eq$ and satisfy the PA recursion equations. The already established
induction supplies every translated induction instance.

All the auxiliary definitions are fixed finite formulas. Once the
interpretation is fixed, its proof translation requires only structural
formula translation, fresh-name choices, finite proof concatenation, and
insertion into the induction template. The supplied paper explicitly gives
this PA-verifiable uniformity. We reviewed that mathematical argument and
replayed its delivered diagnostics, but have not produced a new kernel
formalization of it. Finite diagnostic success alone is not the reason for
Proposition~\ref{prop:theories}.

\section{Native formula codes and the new axiom emitter}

\subsection{The retained codec}

Use dual Cantor pairing and the four constructors
\begin{align*}
 P(x,y)&=\frac{(x+y)(x+y+1)}2+y,\\
 I(x,y)&=4P(x,y),& R(i,j)&=4P(i,j)+1,\\
 Q(i,f)&=4P(i,f)+2,& E_{=}(i,j)&=4P(i,j)+3.
\end{align*}
They represent implication, membership, universal quantification, and
auxiliary equality, respectively. Set $F=1$ and $N(f)=I(f,1)$.
The natural number~0 is a deliberate truth alias, since $I(0,0)=0$.
Every positive number is decoded by its residue modulo~4 and inverse pairing;
every recursive formula field is smaller than the parent. Consequently
all natural numbers are admissible formula codes.

There are two standard equivalent ways to handle falsity. Decode $F$ as
$E(v_0,v_0)$ and normalize it under $\IR$, or decode an implication with
consequent~1 as ordinary negation. In either presentation, inside
$\ThW+\IR$ one has uniform proof transformations for
\[
 D(I(f,g))\leftrightarrow(D(f)\to D(g)),\qquad
 D(Q(i,f))\leftrightarrow\forall v_iD(f).
\]
The zero collapse is harmless because both $\top$ and $\top\to\top$ are
theorems. Treating code~1 as false without justifying irreflexivity would
be a genuine soundness error.

\begin{lemma}[Self-fresh labels]\label{lem:fresh}
For every positive $n$, every variable label in the decoding
of $n$ is strictly smaller than $n$.
\end{lemma}
\begin{proof}
Both coordinates of $P(i,j)$ are at most $P(i,j)$. For positive constructor
codes, the factor~4 and tag therefore make every coordinate strictly
smaller than the parent. Choose the decoding of zero to be a fixed closed
tautology using only label~0. The decoding of~1 is $E(v_0,v_0)$, whose
only label is $0<1$. Strong induction on positive $n$ now proves the claim
for all recursive formula fields, including occurrences of the zero alias.
\end{proof}

\subsection{Fixed source and test variables}

The old axiom selector is $c$. Put $n=c$, using this same number as both
the matrix code and the output-variable label. Define
\begin{align}
 e_n&=R(1,n),\nonumber\\
 G_n&=Q(n,I(R(n,1),R(n,0))),\nonumber\\
 U_n&=Q(1,I(e_n,n)),\label{eq:emitter}\\
 L_n&=Q(1,I(G_n,I(n,e_n))),\nonumber\\
 H_n&=Q(n,\Hbody),\nonumber\\
 V(a,b,c)&=I(R(a,b),Q(c,R(a,b))),\nonumber\\
 B(a,b,c)&=N(I(V(a,b,c),H_c)).\nonumber
\end{align}
The selected table uses \ClauseOrder\ clause order. The two possible orders
differ only by exchanging the conjuncts after classical normalization.
The axiom branch emits $B(a,b,c)$ if $c\ne a$ and $c\ne b$, and emits
truth~0 otherwise. The two old gate tests are retained unchanged.

For $n\ge2$, $G_n$ expresses $\Sub(v_1,v_0)$. Its binder $n$ is local
to that guard and safely shadows the output binder there. The matrix does
not occur inside this guard. By Lemma~\ref{lem:fresh}, $n$ is fresh in
the matrix and differs from the fixed labels~0 and~1. Thus
\begin{equation}\label{eq:shell}
 \neg D(H_n)\ \equiv\ \exists v_n\left[
 \forall v_1(E(v_1,v_n)\to D(n))\ \land\
 \forall v_1(\Sub(v_1,v_0)\to(D(n)\to E(v_1,v_n)))
 \right].
\end{equation}
This is an unrestricted weak-cover instance. The other component of
$D(B)$ is membership vacuity, valid because the original gate makes its
binder fresh. The vacuity component is independent of the fixed source
and tested variables of the power-cover component.

\begin{lemma}[The two exceptional codes]\label{lem:collision}
The outputs corresponding to $n=0$ and $n=1$ are sound in $\ThW+\IR$.
\end{lemma}
\begin{proof}
At $n=0$, the matrix is truth. The binder of $G_0$ captures all
occurrences of label~0 in its own scope, so under $\IR$ it is
\[
 G_0=\forall v_0(E(v_0,v_1)\to E(v_0,v_0))
       \ \leftrightarrow\ \Emp(v_1).
\]
Here $\Emp(x):=\forall t\,\neg E(t,x)$. The upper clause is tautological,
and the shell says
\[
 \exists b\,\forall x\,(\Emp(x)\to E(x,b)).
\]
This follows from $\ThW$. First obtain an empty object $e$ by applying
the false-matrix cover at any source (the domain is nonempty). Then apply
the true-matrix cover at $e$. Since $\Sub(x,e)\leftrightarrow\Emp(x)$,
the resulting object contains every empty object, as required.

At $n=1$, the matrix is $F=E(v_0,v_0)$, false under $\IR$. The upper clause
$U_1$ is $\forall v_1(E(v_1,v_1)\to F)$ and hence expresses
irreflexivity under this normalization. The lower clause holds because
its matrix antecedent is false. All occurrences of label~1 in those
clauses are already bound, so the outer output quantifier is vacuous.
The resulting shell is true. No malformed-input case has been silently
discarded.
\end{proof}

\paragraph{Why no increment or syntax test is needed.}
The local binder of the subset guard is essential to the zero-code
calculation. Replacing that binder by a separately fixed fresh label would
produce a different exceptional formula and would require a new analysis.
For the displayed dynamic guard, both tiny cases are sound and the selector
can be used directly. No run-time matrix traversal or free-variable
validator is required: self-freshness does that work in the proof.

\paragraph{A second pitfall avoided.}
Using arbitrary source and tested labels $b,a$ would require a check
$a\ne b$: when $a=b$, the subset guard is tautological and the schema
becomes unrestricted comprehension. But applying that extra check to the
whole bundle would also suppress diagonal membership vacuity. Fixing the
power-cover roles to distinct labels~0 and~1 avoids both problems and
preserves the old logical interface.

\section{The numerical calculus and its consistency calibration}

Let $\CU$ be the numerical calculus containing truth~0, numerical modus
ponens and generalization, the following eight logical rows, and the gated
bundle~$B$. We display rows as formulas, with bold arguments denoting
formula-code slots and plain arguments denoting variable labels:
\begin{center}\small
\begin{tabular}{cl}
\toprule
Row & Formula shape\\\midrule
$\rho_1$ & $\exists a\,(a=b)$\\
$\rho_2$ & $\forall a(\mathbf b\to\mathbf c)\to
              (\forall a\mathbf b\to\forall a\mathbf c)$\\
$\rho_3$ & $a=b\to(c\in a\to c\in b)$\\
$\rho_4$ & $\forall a\forall b\mathbf c\to\forall b\forall a\mathbf c$\\
$\rho_5$ & $c=a\to(c\in b\to a\in b)$\\
$\rho_6$ & $\exists a\forall a\mathbf c\to\mathbf c$\\
$G$ & $c=a\to(c=b\to\forall_{R(a,b)}(a=b))$\\
$\rho_8$ & $((\mathbf c\to\mathbf u)\to\mathbf a)\to
             ((\mathbf a\to\mathbf c)\to(\mathbf b\to\mathbf c))$\\
\bottomrule
\end{tabular}\end{center}
Existential quantification and negation use the native $F$ abbreviation.
The binder in $G$ is the integer $R(a,b)$, which exceeds both $a$ and $b$.
The actual source emits the dominated family of $\rho_8$, with $u\le a$.
The inherited finite dominated-row derivation recovers unrestricted
$\rho_8$ using only numerical modus ponens, truth, and the elementary
bounds $x,y\le I(x,y)$; it is independent of Theory~K \cite{dominated}.

\begin{lemma}[Retained logical completeness]\label{lem:logic}
The finite equality-free first-order proof templates used by the K215
calculus remain valid when its nonlogical bundle is replaced by $B$.
\end{lemma}
\begin{proof}
The unrestricted $\rho_8$ basis gives classical implication. The row
$\rho_6$, together with weakening, derives $F\to A$ for every code $A$.
Indeed $F\to\exists i\forall i A$ is a weakening instance, and
$\exists i\forall i A\to A$ is $\rho_6$.

The merged row $G$ recovers equality reflexivity, its fresh-binder helper,
and the ordinary Euclidean equality row using only $G$, $\rho_1$,
$\rho_2$, $\rho_6$, modus ponens, generalization, and propositional
logic. The finite derivations are in the retained \texttt{G222} templates.
The new bundle yields every membership-vacuity instance by the classical
projection $\neg(V\to H)\to V$, including the case $a=b$.

One must not infer arbitrary primitive-equality vacuity from this alone.
The inherited proof instead interprets equality by two-sided membership
indiscernibility,
\[
 \operatorname{Eind}(x,y):=
 [\forall r(E(r,x)\leftrightarrow E(r,y))]\land
 [\forall r(E(x,r)\leftrightarrow E(y,r))],
\]
with fresh bound names. The \texttt{L224E}/\texttt{Interp222}
templates establish first-order completeness through that interpretation.
Their premises are exactly the retained logical rows and membership
vacuity just recovered. They do not use the Theory~K schema. Equality-free
formulas, including the $\ThU$ and $\ThW$ schemas, are unchanged by this
interpretation. All templates are finite syntactic proof transformations.
\end{proof}

\begin{lemma}[All needed power-cover instances]\label{lem:cover-complete}
Every ordinary equality-free $\ThW+\IR$ axiom has a uniformly constructible
numerical proof in the retained interpretation.
\end{lemma}
\begin{proof}
Write $\tau$ for quotation, with
$\tau(\neg A)=I(\tau(A),1)$. Numerical identity $I(1,1)$ is both a
propositional theorem and $\tau(\neg E(v_0,v_0))$. Generalization at~0
therefore yields the quotation of $\IR$.

For a desired cover matrix, rename its source and tested variables to~0
and~1, using temporary fresh labels and capture-avoiding bound-variable
renaming. Double-negate the matrix before quotation. This preserves
meaning and guarantees a code $n>1$. Self-freshness then makes $n$ a
legal output label. Choose $c=n$ and $a=b=n+1$. The old gate succeeds,
and projection of $B(a,b,c)$ gives $N(H_n)$. Equation~\eqref{eq:shell}
supplies the desired cover. Remove the double negation by the logical
templates, then universally close and restore the original names by
instantiation and alpha conversion.

For a concrete check on the tiny codes, a true matrix may be padded to
$I(1,1)=16$ and a false one to $I(0,1)=8$, with the meanings understood
under $\IR$. These are proof-compiler
operations; the Turing machine never computes a fresh-variable scan.
\end{proof}

\begin{proposition}[Calibration of the numerical target]\label{prop:calibration}
PA proves
\[
  (\exists p\,\operatorname{Prf}_{\CU}(p,1))
           \ \longleftrightarrow\ \neg\Con(\ThU).
\]
\end{proposition}
\begin{proof}
For soundness, decode a numerical proof into an ordinary proof under
irreflexivity. Every logical row is sound, a failed gate produces truth,
and every successful bundle is a valid unrestricted cover together with
vacuity, including both exceptional cases from
Lemma~\ref{lem:collision}. Apply the bit-membership interpretation in PA.
Code~1 becomes a false diagonal membership assertion, yielding a
PA-refutation, and the supplied lower interpretation translates that
refutation to $\ThU$.

For completeness, a $\ThU$-refutation is a $\ThW+\IR$-refutation.
Lemmas~\ref{lem:logic} and~\ref{lem:cover-complete} translate its axioms
and every inference into the numerical calculus. The standard finite
natural-deduction-to-Tarski and quotation templates turn the contradiction
into code~1. For example, quote a negation of a closed logical truth and
apply numerical modus ponens to that truth. Intermediate auxiliary
equalities are eliminated by membership indiscernibility.

All these maps are primitive recursive on the supplied finite syntax.
Their correctness is verified by PA induction on formulas and proof
lengths. No inference from model-theoretic equiconsistency alone is used.
\end{proof}

\section{Fair proof search and its exact counter program}

The nine logical registers retain the K215 order
\[
 \begin{aligned}
 (N,T,W)&=(\text{next},\text{current},\text{stack}),\\
 (C,D,H)&=(\text{selector},\text{counter},\text{cache}),\\
 (X,Y,L)&=(\text{scratch}_1,\text{scratch}_2,\text{remaining fields}).
 \end{aligned}
\]
At a command boundary the selector and two scratch registers are zero.
The semantic state is $(pl,np,t,w,d,h)=(L,N,T,W,D,H)$.
The stack is nested $P(\text{top},\text{rest})$, with the infinite
zero tail of $\operatorname{unP}(0)=(0,0)$.

If $pl=0$, the next candidate is read from $np$ and $np$ is incremented.
Otherwise fields are read from $pl-1$. Repeated inverse pairing reads
the selector and the indicated additional fields. The commands are:
\begin{description}[style=nextline,leftmargin=1.2em]
\item[{NIP, fields $[0]$.}]
Halt if $h=1$. Otherwise pop $(x,w')=\operatorname{unP}(w)$ and return
$(t,w',d,x)$.
\item[{DROP, fields $[1]$.}]
Return $(h,P(t,w),d,h)$.
\item[{MP$(x)$, fields $[2,x]$.}]
Put $z=d+x$. If $I(t,z)=h$, return $(z,P(0,w),z,0)$;
otherwise return $(0,w,z,0)$.
\item[{AX$(c;x,y)$, fields $[c+3,y,x]$.}]
Put $u=d$, $a=d+x$, and $b=h+y$. Push, in order,
\[
 t,\ Q(u,h),\ \rho_8(a,b,c,u),\ G(a,b,c),\ \rho_6(a,c),
\]
\[
 \rho_5(a,b,c),\ \rho_4(a,b,c),\ \rho_3(a,b,c),\
 \rho_2(a,b,c),\ \rho_1(a,b).
\]
Set current to the gated bundle and set $d=h=0$.
\end{description}
The field order in AX is deliberately $y,x$, whereas the displayed
parameter order is $x,y$. The old counter is $u$, so the directly
generated $\rho_8$ satisfies $u\le a$. AX still pushes exactly ten
stack cells. In particular the change has not silently altered the
proof-compiler stack discipline.

\subsection{Soundness}

The invariant says that $t$, $h$, and every stack entry are $\CU$ theorems;
it imposes no condition on the counter or unread fields. It holds when all
three formula registers are zero. NIP and DROP move theorems; successful
MP applies its tested numerical implication; failed MP supplies truth;
AX supplies the old theorem, its generalization, the logical rows, and
a permitted bundle or truth. A halt therefore extracts a finite numerical
proof of~1 by induction on the command run.

\subsection{Exact compiler macros and completeness}

Write $\mathsf h$ for NIP and $\mathsf d$ for DROP. A block
$\operatorname{Build}(A)$ has the following contract: from every
$(T,w,0,0)$, it halts early or returns $(A,w,0,0)$. Define
\begin{align*}
 \operatorname{Discard}&=\operatorname{MP}(1);\operatorname{AX}(1;0,0);
                             \mathsf h^{10},\\
 J(0)&=\varepsilon,\qquad J(u)=\operatorname{MP}(u)\quad(u>0).
\end{align*}
Discard gives $(0,w,0,0)$: the MP must fail, and the next AX has $c=a=1$,
so the gate fails. From that zero current/cache state, $J(u)$ sets the
counter to $u$ while retaining zero current/cache.
The truth leaf is $\operatorname{Build}(0)=\operatorname{Discard}$.

The following finite blocks compile every proof constructor:
\begin{align*}
 \operatorname{Leaf}_j(a,b,c,u)
 &=\operatorname{Discard};J(u);\operatorname{AX}(c;a-u,b);
      \mathsf h^j;\mathsf d;\mathsf h^{11-j},\\
 \operatorname{BundleLeaf}(a,b,c)
 &=\operatorname{Discard};\operatorname{AX}(c;a,b);\mathsf h^{10},\\
 \operatorname{Build}(Q(i,A))
 &=\operatorname{Build}(A);\mathsf d;J(i);\mathsf h;
   \operatorname{AX}(i;0,0);\mathsf h^9;\mathsf d;\mathsf h;\mathsf h.
\end{align*}
Here $1\le j\le8$ and $u\le a$; for rows other than $\rho_8$ use $u=0$.
BundleLeaf is used for successful gates, $c\ne a,b$.
The finite dominated-row derivation handles an unrestricted $\rho_8$
leaf. For a modus-ponens node with premises $A$ and $I(A,B)$, use
\begin{align*}
 &\operatorname{Build}(A);\mathsf d;\operatorname{Build}(I(A,B));
 \mathsf h;\mathsf d;\mathsf h;\operatorname{MP}(B);\\
 &\hspace{12mm}\operatorname{AX}(B;0,0);\mathsf h^{10};
 \mathsf d;\mathsf h;\mathsf h.
\end{align*}
Substitution in the four exact command equations proves these contracts.
The temporary AX instances in the cleanup have failed gates, and the ten
pops restore the old stack. Encountering cache~1 earlier only helps
completeness. Induction on a finite numerical derivation consequently
constructs a block for its last line.

The actual initialization starts $np$ at~4. This loses no proof.
Let $\operatorname{START}(Y)=\operatorname{AX}(0;1,Y)$, which resets
$d=h=0$ from any retained tuple. A compiled proof of~1 gives a command
list
\[
 \operatorname{START}(Y);\operatorname{Build}(1);
                   \mathsf d;\mathsf h;\mathsf h
\]
that halts from every retained tuple. Encode a noninitial command by
$1+$ its nested pairing of fields, and encode the first command without
that extra~1. If $s$ is the encoded suffix after START, the candidate is
\[
 N_Y=P(3,P(Y,P(1,s)))\ge Y.
\]
Choose $Y=4$. At positive $pl$, every completed command leaves a remainder
at most $pl-1$, since each inverse-pairing coordinate is bounded by its
input. Thus every candidate drains. Induction on $N_Y-4$ proves that
the enumerator either halts earlier or eventually starts candidate $N_Y$.
This is a cofinal, constructive argument, valid for internally finite
nonstandard codes in PA as well as standard ones.

\begin{proposition}\label{prop:enum}
PA proves that the counter source from $(N,T,W,C,D,H,X,Y,L)
=(4,0,0,0,0,0,0,0,0)$ halts if and only if a finite $\CU$ proof of~1 exists.
\end{proposition}

\subsection{The exact source replacement}

The accepted K215 source has 4,751 primitive occurrences. We replace its
closed interval $[2885,4482)$, the old Theory~K shell, by a closed expression
with 737 primitive occurrences. The result has 3,891 occurrences. The
other 3,154 primitive instructions and every associated control-flow edge
are independently checked after the ordinal index translation; no retained
edge enters the interior of the replacement.

The inherited arithmetic macros satisfy, for arbitrary natural registers
and cleared scratch values,
\begin{align*}
 \operatorname{PUSH0}:&(T,W)\mapsto(0,P(T,W)),\\
 \operatorname{PUSH1}:&(T,W)\mapsto(1,P(T,W)),\\
 \operatorname{PAR3}:&(T,W)\mapsto(C,P(T,W)),\\
 K:&(T,P(A,W))\mapsto(K(A,T),W)\quad(K\in\{I,R,Q\}),\\
 N:&T\mapsto I(T,1).
\end{align*}
The third line preserves $C$; every macro restores both scratch registers
to~0. The new atom for $n$ is exactly PAR3.
Structural induction on its postfix expression proves that the new block
computes exactly $H_C$ and pushes the old current once. At this point
the old current is $V(a,b,c)$; the unchanged final implication and negation
consume it and produce $B(a,b,c)$, with no additional net stack cell.

The pair, unpair, copy, and destructive-comparison loop contracts are the
retained all-natural counter-program lemmas, including
\texttt{cons119}, \texttt{par168}, \texttt{pushwff171}, \texttt{v177},
\texttt{wal180}, \texttt{wel181}, and \texttt{wn218} in the included
\texttt{ProgLibPerm.v}. An independent check binds their source DAG
expansions to the exact reused intervals: resolved control edges agree,
with only adjacent commuting increments reordered and no control edge
entering the interior of such a pair. The pair/unpair loop lemmas are
\texttt{pair162} and \texttt{unpair115}/\texttt{unpair116}.
Their termination is proved by decreasing natural-number variants.
The new work is a finite
composition of those same closed macros. Parenthesization changes only
closed sequencing, and both independent source-CFG resolvers verify the
control flow. Hence each abstract command has a finite source execution
with exactly the displayed effect, and the only halting instruction remains
the NIP target test. This proves the source part of
Proposition~\ref{prop:enum}; literal-table simulation is the next obligation.

\section{Verification of the literal binary machine}

\subsection{Canonical complete tape frames}

A source boundary frame has a left sentinel~1, a binary program address
$p$, enough zero padding to make the address area length $64+g$, a
header~1, nine primary unary register blocks, and a nonempty finite list
of spare unary blocks. Both exterior half-lines are zero.
A register value $v$ is represented by
\[
 B(v)=0\,1^{v+1}.
\]
The logical-to-physical register permutation is
$(6,2,0,4,3,7,8,5,1)$. The extra gap $g$, all nine values, the first
spare value, and the remaining finite spare list are arbitrary.
An increment appends a spare block of value~1; a successful decrement
increases the gap by one. Address overflow moves the left sentinel one
cell left and increases the gap correspondingly. These effects are
included in the target frame, not discarded as irrelevant tape garbage.

\subsection{Three sound symbolic proof rules}

The checker in \texttt{code/runtime} represents unary lengths by affine
expressions in independent natural variables. It uses only:
\begin{enumerate}
\item a literal transition when the scanned run has a guaranteed positive
constant term;
\item induction across an arbitrary unary run when the exact table row
is a self-loop moving in the scan direction;
\item induction across an arbitrary finite list of blocks after separately
proving the one-block transformation with an arbitrary value and untouched
opaque contexts.
\end{enumerate}
The zero-length case of a unary scan is the identity. An unknown length
cannot be consumed by a literal transition. Reversing the head direction
reverses the represented tape-side orientation as well. Empty spare tails
are allowed, while the separate first spare ensures a nonempty whole list.

The only additional block family is $C(v)=1\,0\,1^v$.
There are four one-block schemas, each taking exactly $v+2$ steps:
\begin{center}
\begin{tabular}{ccc}
\toprule
Native entry state & Direction & Block transformation\\\midrule
64 & right & $B(v)\mapsto B(v)$\\
0 & right & $B(v)\mapsto C(v)$\\
40 & left & $B(v)\mapsto B(v)$\\
97 & left & $C(v)\mapsto B(v)$\\
\bottomrule
\end{tabular}
\end{center}
For example, the first uses the actual rows
$64,0\mapsto0R74$, $74,1\mapsto1R64$, and then the
$64,1\mapsto1R64$ self-loop $v$ times. The full four derivations are
recorded in the runtime report and independently reviewed.

The normalization identity $C(A)1=1B(A)$ for a finite block list~$A$,
and its reversal, follow by list induction. The tail-width variable means
$\sum_{v\in A}(v+2)$, not the number of blocks. It is preserved by the
block transformations. Allowing it to be an independent natural variable
is a sound strengthening of the checked claim.

For every source primitive and, for a decrement, for both zero and positive
inputs, the checker starts from the full symbolic frame and reaches the
specified successor frame or the halt state. It also requires a duration
with nonnegative coefficients and a strictly positive constant term. The
result therefore supports finite concatenation and rules out infinitely
many source steps occurring in only finitely many machine steps.
The limit on symbolic proof construction bounds the checker's search for a
derivation; it does not bound register values or actual run times.

\begin{lemma}[Literal simulation]\label{lem:runtime}
For the released table and source, every canonical primitive frame makes
the corresponding source transition after a positive finite number of
literal machine steps, with the exact complete successor tape frame.
\end{lemma}
\begin{proof}
The complete symbolic run succeeds for all \Ncases\ cases. Soundness of
each rule follows from the preceding unary or finite-list induction, and
the only noninductive rule is an exact table lookup. Finite composition
proves each case; the complete source census proves coverage. The checks
use every value as an arbitrary natural parameter, rather than sampling
large register values. All rule soundness arguments and the finite
compositions are expressible and provable in PA.
\end{proof}

\subsection{Blank startup and the final composition}

Starting with no 1-cells, the literal machine reaches the root dispatcher
after \BootSteps\ transitions, with head at~$-8$, left sentinel at~$-9$,
header at~57, and extra gap~1. Its primary logical registers are
$(4,0,0,0,0,0,0,0,0)$, followed by forty spare blocks of value~1. There
are exactly 95 one-cells. An independent simulator checks every transition
of this finite prefix and the complete final finite support; the complement
of that support is zero. The startup theorem is a separate condition from
the arbitrary-register simulation report.

The source-to-enumerator proof and Lemma~\ref{lem:runtime} give both
directions of halting equivalence. There is no unaccounted halt in a
nonhalting primitive simulation: crossing the halt state would make the
next literal lookup fail instead of proving the required successor frame.
Combining startup, Proposition~\ref{prop:enum}, and
Proposition~\ref{prop:calibration} gives a PA proof of
\[
 \Halt(M_U,\mathrm{blank})\longleftrightarrow\neg\Con(\ThU).
\]
Classical propositional logic yields Theorem~\ref{thm:main}.

\section{Construction history, reproducibility, and limits}

\subsection{Why the physical checks matter}

The first source layout produced a 228-state table that passed decoder
prefix checks and exact blank startup, but failed a primitive simulation.
A reused increment entry selected its continuation using the preceding
tape bit. Its new location did not have the required bit. A concrete
all-zero-register execution reproduced the discrepancy. A later 222-state
prefix allocation had the same class of defect. These rejected tables
are not claimed as results.

The corrected allocation checks the contexts of retained arithmetic
entries and protects the intermediate rows used by their marker-dependent
routing. Every delivered table is then subjected to the complete
symbolic runtime check. The final result depends on that check, not on
the allocator's estimate or on the count of distinct source subtrees.
\begin{center}
\begin{tabular}{lrl}
\toprule
Layout & States & Decisive verification status\\\midrule
Initial shifted emitter & 228 & Rejected: wrong arithmetic context\\
Early shared shifted emitter & 222 & Rejected: wrong arithmetic context\\
Corrected shifted emitter & 224 & Complete runtime and startup pass\\
Further shared shifted emitter & 223 & Complete runtime and startup pass\\
Unshifted dynamic-guard emitter & 222 & Complete runtime and startup pass\\
\bottomrule
\end{tabular}
\end{center}
The selected 222-state table is a different literal machine from the
rejected earlier table of the same size. Its exact bytes, not its short
name or state count, identify the result. The selected allocation removes
one old working state and adds eight: $215-1+8=222$. It changes 72 old
transition rows and adds 16 fresh rows. Strong transition-bisimulation
refinement leaves every working state in a separate class; this rules out
that immediate merge operation, not more general smaller constructions.


\Needspace{9\baselineskip}
\subsection{Reproduction}

The archive is self-contained for reconstruction and verification of the
new table. It needs Python~3.10 or later and the standard library. No
solver, network access, or source repository checkout is needed. Run:
\begin{lstlisting}
python3 reproduce.py
\end{lstlisting}
This checks release hashes, reconstructs the selected source from the
frozen K215 input and explicit sequencing recipe, reconstructs both
transition tables from the exact finite allocation recipe, runs the
independent source audit, verifies the state-renaming bijection, and
replays initialization. To repeat every arbitrary-register symbolic case,
run:
\begin{lstlisting}
python3 reproduce.py --full
\end{lstlisting}
The full run also requires the successful startup condition and checks
the complete expected case count. It may take several minutes.
The saved runtime case file contains reproducible summaries, not a
standalone proof object with every symbolic event. The small checker,
table, and source regenerate the full verification.

The original supplied $\ThU$ report and its own reproduction code are
included intact. From \texttt{references/theory\_U}, its command is
\texttt{python3 -B reproduce.py}. Its six diagnostic suites run in normal
and optimized Python modes; all sixteen named certificate targets and
the 4,096-description bounded searches were replayed. These finite
checks neither prove consistency nor replace its interpretation proof.

The LaTeX source is \texttt{paper/report.tex}, with exact machine metadata
and the full table appendix generated from the delivered data. Compile
from that directory with \texttt{pdflatex report.tex} twice. The released
PDF was rendered and inspected; different TeX installations and timestamps
may change PDF bytes without changing its mathematical content.

\subsection{Formal status}

The proof consists of ordinary mathematical arguments about finite syntax,
natural-number counter programs, and literal binary transitions. Each
unbounded step uses induction available in PA. The new symbolic checker
has been independently reviewed at the level of its proof rules. It is
not a proof-assistant kernel, and this release supplies no new checked
MM0 theorem with the final literal table in its statement. The existing
K215 MM0 result is relevant evidence for the reusable architecture, but
does not certify this modification automatically.

If $\PA$ is consistent, Theorem~\ref{thm:main} implies actual non-halting
and PA-unprovability of this non-halting assertion by the second
incompleteness theorem. Mere consistency does not by itself justify the
stronger statement that PA also cannot prove halting; that assertion needs
an appropriate soundness hypothesis. Our main theorem makes neither
extra assertion part of its proof.

\subsection{Why proof-theoretic strength does not order state counts}

A consistency-search program's size depends on its axiom generator,
syntax, proof engine, and compiler. Proof-theoretic strength places no
monotone lower bound on those sizes: a stronger theory can have a shorter
presentation. Here the arithmetic interpretation appears only in the
correctness proof. The machine emits a short covering schema. Its state
cost is determined by using formula codes as fresh variable labels,
fixing source and test roles, retaining the membership-vacuity interface,
and sharing physical states.

\appendix
\section{Exact transition table}

Each entry is written $wDq$, meaning write $w$, move $D\in\{L,R\}$,
and enter state $q$. The start is~0; the halt label is~\Nstates.
The table's SHA-256, in the supplied ASCII file format, is
\begin{center}\small\texttt{\DenseHash}\end{center}
The native table SHA-256 is
\begin{center}\small\texttt{\NativeHash}\end{center}
The source-site SHA-256 is
\begin{center}\small\texttt{\SourceHash}\end{center}
% Generated from every row of machine/u.tm.
\small
\begin{longtable}{@{}rll@{\hspace{1.2em}}rll@{\hspace{1.2em}}rll@{}}
\toprule
$q$ & read 0 & read 1 & $q$ & read 0 & read 1 & $q$ & read 0 & read 1\\
\midrule\endfirsthead
\toprule
$q$ & read 0 & read 1 & $q$ & read 0 & read 1 & $q$ & read 0 & read 1\\
\midrule\endhead
\bottomrule\endfoot
0 & \texttt{1R93} & \texttt{1R0} & 74 & \texttt{1R0} & \texttt{1R64} & 148 & \texttt{0R85} & \texttt{1R89}\\
1 & \texttt{0L1} & \texttt{1L2} & 75 & \texttt{0R75} & \texttt{1R83} & 149 & \texttt{0R161} & \texttt{1R12}\\
2 & \texttt{0L3} & \texttt{1L2} & 76 & \texttt{0L147} & \texttt{1L143} & 150 & \texttt{0R84} & \texttt{1R145}\\
3 & \texttt{0L5} & \texttt{1L2} & 77 & \texttt{0R86} & \texttt{1R87} & 151 & \texttt{0R217} & \texttt{1R130}\\
4 & \texttt{1R6} & \texttt{0L4} & 78 & \texttt{0R11} & \texttt{1R131} & 152 & \texttt{0L67} & \texttt{1L143}\\
5 & \texttt{0L9} & \texttt{1L2} & 79 & \texttt{0R90} & \texttt{1L134} & 153 & \texttt{0R198} & \texttt{1R151}\\
6 & \texttt{0R6} & \texttt{0R10} & 80 & \texttt{0L109} & \texttt{1R57} & 154 & \texttt{0R161} & \texttt{1R116}\\
7 & \texttt{0R85} & \texttt{1R194} & 81 & \texttt{0R49} & \texttt{0R33} & 155 & \texttt{0R98} & \texttt{1R98}\\
8 & \texttt{0R14} & \texttt{1R13} & 82 & \texttt{0R92} & \texttt{0L42} & 156 & \texttt{0R49} & \texttt{0R32}\\
9 & \texttt{0L15} & \texttt{1L2} & 83 & \texttt{0R94} & \texttt{1R162} & 157 & \texttt{0R88} & \texttt{1R113}\\
10 & \texttt{0L16} & \texttt{0R10} & 84 & \texttt{0R97} & \texttt{1R12} & 158 & \texttt{0R193} & \texttt{1R171}\\
11 & \texttt{0L26} & \texttt{1R115} & 85 & \texttt{0R89} & \texttt{1R98} & 159 & \texttt{0L192} & \texttt{1R97}\\
12 & \texttt{0R19} & \texttt{1R20} & 86 & \texttt{0R212} & \texttt{1R21} & 160 & \texttt{0R170} & \texttt{1R183}\\
13 & \texttt{0L90} & \texttt{1L31} & 87 & \texttt{0R118} & \texttt{1R101} & 161 & \texttt{0R19} & \texttt{1R173}\\
14 & \texttt{0L11} & \texttt{0R22} & 88 & \texttt{0R220} & \texttt{1R153} & 162 & \texttt{0R109} & \texttt{1R26}\\
15 & \texttt{0L23} & \texttt{1L2} & 89 & \texttt{0R103} & \texttt{1R104} & 163 & \texttt{0R161} & \texttt{1R97}\\
16 & \texttt{1L1} & \texttt{0L1} & 90 & \texttt{0R81} & \texttt{1R80} & 164 & \texttt{0R127} & \texttt{1R166}\\
17 & \texttt{0R24} & \texttt{0R32} & 91 & \texttt{0R102} & \texttt{1R200} & 165 & \texttt{0R98} & \texttt{1R164}\\
18 & \texttt{0R150} & \texttt{1R178} & 92 & \texttt{0R106} & \texttt{1R107} & 166 & \texttt{0L7} & \texttt{1R144}\\
19 & \texttt{0R28} & \texttt{1R29} & 93 & \texttt{0R66} & \texttt{0R0} & 167 & \texttt{0L110} & \texttt{1R177}\\
20 & \texttt{0R30} & \texttt{1R31} & 94 & \texttt{0R27} & \texttt{1R108} & 168 & \texttt{0R148} & \texttt{1R176}\\
21 & \texttt{0R24} & \texttt{0R60} & 95 & \texttt{0R187} & \texttt{0L213} & 169 & \texttt{0R218} & \texttt{1R204}\\
22 & \texttt{0R33} & \texttt{1R22} & 96 & \texttt{1L111} & \texttt{1L96} & 170 & \texttt{0R163} & \texttt{1R145}\\
23 & \texttt{0L34} & \texttt{1L2} & 97 & \texttt{0R19} & \texttt{1R112} & 171 & \texttt{0R199} & \texttt{1R125}\\
24 & \texttt{0L35} & \texttt{0R0} & 98 & \texttt{0R61} & \texttt{1R61} & 172 & \texttt{0R99} & \texttt{1R219}\\
25 & \texttt{0R36} & \texttt{1R25} & 99 & \texttt{0R69} & \texttt{1R137} & 173 & \texttt{0R185} & \texttt{1R124}\\
26 & \texttt{0R37} & \texttt{1R38} & 100 & \texttt{0R146} & \texttt{1R141} & 174 & \texttt{0R147} & \texttt{1R167}\\
27 & \texttt{0R39} & \texttt{0L40} & 101 & \texttt{0R106} & \texttt{0L213} & 175 & \texttt{0R17} & \texttt{1R182}\\
28 & \texttt{0R41} & \texttt{1R42} & 102 & \texttt{0R119} & \texttt{1R120} & 176 & \texttt{0R76} & \texttt{1R61}\\
29 & \texttt{0R43} & \texttt{1R44} & 103 & \texttt{0R72} & \texttt{1R59} & 177 & \texttt{0R147} & \texttt{1R89}\\
30 & \texttt{0R45} & \texttt{1R46} & 104 & \texttt{0R58} & \texttt{1R121} & 178 & \texttt{0R154} & \texttt{1R206}\\
31 & \texttt{0R21} & \texttt{1R47} & 105 & \texttt{0R21} & \texttt{1R122} & 179 & \texttt{0R142} & \texttt{1R85}\\
32 & \texttt{0R48} & \texttt{1R32} & 106 & \texttt{0R21} & \texttt{0L10} & 180 & \texttt{0R179} & \texttt{1R76}\\
33 & \texttt{0R32} & \texttt{1R33} & 107 & \texttt{0R62} & \texttt{1R42} & 181 & \texttt{0R148} & \texttt{1R216}\\
34 & \texttt{0L50} & \texttt{1L2} & 108 & \texttt{0R123} & \texttt{1R8} & 182 & \texttt{0R78} & \texttt{0L213}\\
35 & \texttt{1R35} & \texttt{1L1} & 109 & \texttt{0R102} & \texttt{1R124} & 183 & \texttt{0R98} & \texttt{1R52}\\
36 & \texttt{0R51} & \texttt{1R36} & 110 & \texttt{0R152} & \texttt{1R191} & 184 & \texttt{0R170} & \texttt{1R52}\\
37 & \texttt{0R17} & \texttt{1R211} & 111 & \texttt{0L213} & \texttt{0L96} & 185 & \texttt{0R17} & \texttt{1L133}\\
38 & \texttt{0R13} & \texttt{1R53} & 112 & \texttt{0R126} & \texttt{1R105} & 186 & \texttt{0R145} & \texttt{1R100}\\
39 & \texttt{0L11} & \texttt{0R36} & 113 & \texttt{0R215} & \texttt{1R218} & 187 & \texttt{0R77} & \texttt{1R118}\\
40 & \texttt{0L4} & \texttt{1L54} & 114 & \texttt{0R167} & \texttt{1R185} & 188 & \texttt{0R107} & \texttt{1R126}\\
41 & \texttt{0R24} & \texttt{0R25} & 115 & \texttt{0R76} & \texttt{1R207} & 189 & \texttt{1L40} & \texttt{1L111}\\
42 & \texttt{0R13} & \texttt{0L16} & 116 & \texttt{0R150} & \texttt{1R148} & 190 & \texttt{0R174} & \texttt{1R161}\\
43 & \texttt{0R56} & \texttt{1R57} & 117 & \texttt{0R67} & \texttt{1R181} & 191 & \texttt{0R152} & \texttt{1R117}\\
44 & \texttt{0R58} & \texttt{1R59} & 118 & \texttt{0L133} & \texttt{0L1} & 192 & \texttt{0R186} & \texttt{1R187}\\
45 & \texttt{0R24} & \texttt{0R55} & 119 & \texttt{0R39} & \texttt{1R73} & 193 & \texttt{0R204} & \texttt{1R161}\\
46 & \texttt{0R61} & \texttt{1R62} & 120 & \texttt{0R62} & \texttt{1R134} & 194 & \texttt{0R159} & \texttt{1R195}\\
47 & \texttt{0R63} & \texttt{0L16} & 121 & \texttt{0R62} & \texttt{1R135} & 195 & \texttt{0R179} & \texttt{1R190}\\
48 & \texttt{0R25} & \texttt{1R48} & 122 & \texttt{0R71} & \texttt{0L16} & 196 & \texttt{0R65} & \texttt{1R7}\\
49 & \texttt{0L35} & \texttt{0R64} & 123 & \texttt{0R14} & \texttt{1R134} & 197 & \texttt{0R152} & \texttt{1R129}\\
50 & \texttt{0L75} & \texttt{1L2} & 124 & \texttt{0R21} & \texttt{1R136} & 198 & \texttt{0R161} & \texttt{1R130}\\
51 & \texttt{0R16} & \texttt{1R51} & 125 & \texttt{0R203} & \texttt{1R161} & 199 & \texttt{0R161} & \texttt{1R204}\\
52 & \texttt{0R149} & \texttt{1R145} & 126 & \texttt{0R139} & \texttt{1L20} & 200 & \texttt{0R105} & \texttt{1R97}\\
53 & \texttt{0R68} & \texttt{1R69} & 127 & \texttt{0R132} & \texttt{1R184} & 201 & \texttt{0R91} & \texttt{1R89}\\
54 & \texttt{0L40} & \texttt{1L54} & 128 & \texttt{0R209} & \texttt{1R174} & 202 & \texttt{0R196} & \texttt{1R117}\\
55 & \texttt{0R22} & \texttt{1R55} & 129 & \texttt{0R76} & \texttt{1R180} & 203 & \texttt{0R89} & \texttt{1R61}\\
56 & \texttt{0L11} & \texttt{0R51} & 130 & \texttt{0R193} & \texttt{1R203} & 204 & \texttt{0R19} & \texttt{1R61}\\
57 & \texttt{0R61} & \texttt{0L16} & 131 & \texttt{0R208} & \texttt{1R159} & 205 & \texttt{0R175} & \texttt{1R12}\\
58 & \texttt{0R70} & \texttt{1R71} & 132 & \texttt{0R12} & \texttt{1R143} & 206 & \texttt{0R221} & \texttt{1R157}\\
59 & \texttt{0R62} & \texttt{1R72} & 133 & \texttt{0R30} & \texttt{1R70} & 207 & \texttt{0R163} & \texttt{1R210}\\
60 & \texttt{0R55} & \texttt{1R60} & 134 & \texttt{0R39} & \texttt{1L126} & 208 & \texttt{0R201} & \texttt{1R28}\\
61 & \texttt{0R49} & \texttt{0R25} & 135 & \texttt{0R155} & \texttt{0L16} & 209 & \texttt{0R165} & \texttt{1R197}\\
62 & \texttt{0L28} & \texttt{1L73} & 136 & \texttt{0R156} & \texttt{0L16} & 210 & \texttt{0R172} & \texttt{1R95}\\
63 & \texttt{0L119} & \texttt{1L30} & 137 & \texttt{0R89} & \texttt{1R138} & 211 & \texttt{0R205} & \texttt{0L213}\\
64 & \texttt{0R74} & \texttt{1R64} & 138 & \texttt{0R201} & \texttt{1R142} & 212 & \texttt{0R45} & \texttt{1R101}\\
65 & \texttt{0R98} & \texttt{1R152} & 139 & \texttt{0R24} & \texttt{0R48} & 213 & \texttt{0L4} & \texttt{0L213}\\
66 & \texttt{1R189} & \texttt{1L54} & 140 & \texttt{0R203} & \texttt{1R218} & 214 & \texttt{0R166} & \texttt{1R192}\\
67 & \texttt{0R146} & \texttt{1R147} & 141 & \texttt{0R202} & \texttt{1R114} & 215 & \texttt{0R161} & \texttt{1R167}\\
68 & \texttt{0R77} & \texttt{1R222} & 142 & \texttt{0R97} & \texttt{1R161} & 216 & \texttt{0R128} & \texttt{1R18}\\
69 & \texttt{0L201} & \texttt{1R31} & 143 & \texttt{0R179} & \texttt{1R168} & 217 & \texttt{0R89} & \texttt{1R169}\\
70 & \texttt{0R41} & \texttt{1R63} & 144 & \texttt{0R12} & \texttt{1R161} & 218 & \texttt{0R89} & \texttt{1R85}\\
71 & \texttt{0L90} & \texttt{1R79} & 145 & \texttt{0R89} & \texttt{1L168} & 219 & \texttt{0R214} & \texttt{1R188}\\
72 & \texttt{0R56} & \texttt{1L80} & 146 & \texttt{0R89} & \texttt{1R160} & 220 & \texttt{0R19} & \texttt{1R140}\\
73 & \texttt{0R81} & \texttt{1R82} & 147 & \texttt{0L154} & \texttt{1L165} & 221 & \texttt{0R161} & \texttt{1R158}\\
\end{longtable}
\normalsize


\section{Pinned sources and attribution}

The supplied $\ThU$ archive has SHA-256
\begin{center}\small
\texttt{0fcb1117a5353766a3366bab04e56b4efbd0607833cdc73ffb34b67e6808b0ff}.
\end{center}
The ProveIt checkout is pinned to
\texttt{651d6258e9d57f974839ff2ef98060466973a512}; the
\texttt{k215-pa} reference to
\texttt{992b76d60d1a138d84e184cd8b6796243646a984}; and the
\texttt{zf-254} reference to
\texttt{c132d7eae8158cd3a9850ea3eb7251158c3dd821}.

The new work is the fixed-role, code-as-label power-covering source,
its logical calibration to $\ThU$, the selected source layout and corrected
physical allocation, and the exact verification and packaging of the
resulting table. The proof-search architecture, native codec, logical
templates, source-control language, and register interpreter are inherited
from the cited Mira, Wade, Jeff Ketchersid, Elliot Glazer, and ProveIt work.
The theory $\ThU$ and its arithmetic interpretation are inherited from the
user-supplied report. File-level provenance and retained notices are
listed in \texttt{NOTICE.txt}; no repository-wide license is newly asserted
for upstream material whose licensing is unspecified.

\begin{thebibliography}{9}
\bibitem{unary} Supplied research package, \emph{Unary Power Covering at PA
Strength}, 27 October 2026. Exact input:
\url{https://github.com/VladimirReshetnikov/ProveIt/blob/main/docs/incoming/pa_unary_power_cover.zip}.
Included under \texttt{references/theory\_U}.
\bibitem{proveit} V. Reshetnikov and contributors, \emph{ProveIt, accepted
K215 construction and verification}, pinned revision above.
\url{https://github.com/VladimirReshetnikov/ProveIt/tree/main/Computability/BusyBeaver/Research/ZFC/candidate-k215}.
\bibitem{k215pa} E. Glazer and contributors, \emph{K215 halts iff ZFC is
inconsistent, formalized in Metamath Zero}.
\url{https://github.com/elliotglazer/k215-pa}.
\bibitem{k230} E. Glazer and contributors, \emph{K230: a 230-state Turing
machine for the consistency of ZFC}, especially the Theory~K description.
\url{https://github.com/elliotglazer/k230-rocq}.
\bibitem{m254} J. Ketchersid and contributors, \emph{Mira's 254-state
construction: logic, program, and binary-machine verification}.
\url{https://github.com/jeff87654/busy-beaver/tree/zf-254/ZF/m254}.
\bibitem{m238} Mira and verification contributors, \emph{The 238-state
PA-consistency machine}, upstream \texttt{PA/m238}.
\url{https://github.com/jeff87654/busy-beaver/tree/main/PA/m238}.
\bibitem{dominated} ProveIt finite logical templates, particularly
\texttt{DominatedCore.v}, \texttt{DomRho8.v}, \texttt{G222.v},
\texttt{L224E.v}, \texttt{Interp222.v}, and \texttt{Complete222.v}, in
\texttt{k225-formal/rocq/rocq-extension/coq/215}; arithmetic counterparts
\texttt{k215\_rho8\_recovery.mm0} and
\texttt{k218\_translation\_logic.mm0} in \cite{k215pa}.
\end{thebibliography}
\end{document}
